% !TEX program = tectonic
\documentclass[9pt,a4paper]{extarticle}
\usepackage[top=0.45cm,bottom=0.38cm,left=0.85cm,right=0.85cm]{geometry}
\usepackage[T1]{fontenc}\usepackage[utf8]{inputenc}\usepackage{lmodern}
\usepackage[french]{babel}
\usepackage{amsmath,amssymb}\usepackage{textcomp}\usepackage{xcolor}\usepackage{booktabs}\usepackage{array}\usepackage{multirow}
\newcolumntype{R}[1]{>{\raggedright\arraybackslash}p{#1}}
\usepackage{enumitem}\usepackage{multicol}\usepackage{graphicx}
\usepackage[hidelinks]{hyperref}\usepackage{url}\urlstyle{sf}
\graphicspath{{figs_nb13/}{figs_nb15/}}
\definecolor{bleu}{RGB}{20,77,144}\definecolor{vert}{RGB}{20,110,60}
\definecolor{rouge}{RGB}{160,30,30}\definecolor{grisf}{RGB}{110,110,110}
\definecolor{or}{RGB}{150,105,0}
\newcommand{\fig}[2]{\begin{center}\includegraphics[width=0.86\linewidth]{#1}\end{center}\vspace{-2.5mm}{\scriptsize\color{grisf} #2}\par\vspace{0.6mm}}
\pagestyle{empty}\setlength{\parindent}{0pt}\setlength{\parskip}{0.45pt}
\setlength{\abovedisplayskip}{2.5pt}\setlength{\belowdisplayskip}{2.5pt}
\setlength{\abovedisplayshortskip}{1pt}\setlength{\belowdisplayshortskip}{1pt}
\renewcommand{\arraystretch}{0.95}
\setlength{\columnsep}{5mm}\setlength{\columnseprule}{0.3pt}
\renewcommand{\columnseprulecolor}{\color{grisf!40}}
% un bloc = le niveau de titre fort (double filet)
\newcommand{\mouv}[2]{\vspace{1.7mm}{\color{bleu}\hrule height 1pt}\vspace{0.8mm}%
  {\color{bleu}\bfseries #1~---~#2}\par\vspace{0.5mm}{\color{bleu!35}\hrule height 0.5pt}\vspace{1.1mm}}
% un sous-titre à l'intérieur d'un bloc
\newcommand{\et}[2]{\vspace{1.05mm}{\color{bleu}\bfseries #1 $\cdot$ #2}\par\vspace{0.2mm}}
% une étape du protocole
\newcommand{\etape}[2]{\vspace{0.9mm}{\color{rouge}\bfseries #1.}~{\bfseries #2}\par\vspace{0.15mm}}
% ce que le notebook imprime pour attester l'étape
\newcommand{\sortie}[1]{\vspace{0.4mm}{\scriptsize\color{vert}\textbf{\checkmark}~#1}\par\vspace{0.2mm}}
\newcommand{\lecon}[1]{\vspace{0.4mm}{\scriptsize\color{vert}$\blacktriangleright$~\textbf{ce que ça apprend :} #1}\par\vspace{0.2mm}}
\newcommand{\obs}[1]{\vspace{0.2mm}\hangindent=3.4mm\hangafter=0{\color{bleu}\footnotesize\textbullet}~{\small #1}\par\vspace{0.2mm}}
\newenvironment{pts}{\begin{itemize}[leftmargin=3.4mm,itemsep=0.2mm,topsep=0.4mm,parsep=0pt,label={\color{bleu}\footnotesize\textbullet}]}{\end{itemize}}
\begin{document}

\begin{center}
{\large\bfseries M3 --- pondérer au dollar \emph{et} au nombre d'élus}
\end{center}
\vspace{0.4mm}\hrule\vspace{1.2mm}

\begin{multicols}{2}

% ══════════════════════════════════════════════════════════════════════
\mouv{I --- Ce qu'on fait}{dix étapes, dans l'ordre d'exécution}

\etape{1}{Les données}
\begin{pts}
\item une source : la table des déclarations ;
\item trois filtres : classe d'actif = \textbf{action}, transaction entre 2014 et 2026, série de prix
      exploitable ($\ge 2$ cotations, cours $\ge 0{,}10$~\$, aucun saut journalier au-delà de 300~\%) ;
\item chaque opération $k$ porte : l'élu $m(k)$, son parti $P(k)$, le titre $i(k)$, le sens, la date de
      transaction $\tau_k$, celle de \textbf{divulgation} $\delta_k$, le montant $A_k$ = milieu de fourchette ;
\item il reste \textbf{113\,369 opérations}, \textbf{350 élus} et \textbf{2\,755 titres}.
\end{pts}

\etape{2}{Le calendrier et les prix}
\begin{pts}
\item $\mathcal C$ = les 3\,645 jours cotés du SPY (03/01/2012 $\to$ 02/07/2026) ; toute date est un indice ;
\item $P_i(t)$ = cours de clôture, reporté sur $\mathcal C$ par le dernier cours connu ;
\item un titre est \emph{vivant} en $t$ si sa série a commencé et n'est pas arrêtée depuis plus de 4 jours.
\end{pts}

\etape{3}{Les coupes}
$t$ parcourt le dernier jour coté de chaque mois, de janvier 2014 à juin 2026 : 150 coupes. On garde la
première où l'univers atteint $\lceil 1/c\rceil = 10$ titres \textbf{pour les deux partis}, et les suivantes.

\etape{4}{La fenêtre du signal}
À la coupe $t$, parti $P$ : \textbf{les achats seuls}, datés à leur \textbf{divulgation}.
\[ \mathcal K^P(t) = \bigl\{\,k : P(k)=P,\ \text{achat},\ t - W_3 < \delta_k \le t \,\bigr\},
   \quad W_3 = 756\ \text{j} \]
Les ventes n'entrent pas dans le score, ni récentes ni anciennes.

\etape{5}{La clé d'agrégation, puis le score}
La clé est le \textbf{titre}, classes d'actions fusionnées (GOOGL$\to$GOOG, BRK-A$\to$BRK-B, FOXA$\to$FOX,
NWSA$\to$NWS). \emph{C'est elle, et elle seule, que le bloc III remplacera par un ETF.}
\[ s_{i} \;=\; \underbrace{\textstyle\sum_{k} A_k}_{B_i \ \text{: les dollars}}
   \;\times\; \underbrace{\#\{\,m(k)\,\}}_{m_i \ \text{: les élus distincts}} \]
$K$ élus au même montant $a$ donnent $s_i = K^2 a$ : \textbf{le consensus compte au carré}. Les deux comptes
se font après la fusion.

\etape{6}{L'univers}
Parmi les clés vivantes en $t$ et cotées au close $t+1$ :
$\mathcal S(t) = \text{les } N = 150 \text{ plus gros } s_i$.

\etape{7}{Les poids}
Normalisation, puis plafond $c = 10$~\% par \textbf{remplissage de niveau} --- les lignes saturées y restent,
les autres montent jusqu'à ce que la somme fasse 1 :
\[ \boxed{\;w_i \;=\; \min\bigl(c,\; \lambda\, v_i\bigr), \qquad v_i = s_i \Big/ \textstyle\sum_{j\in\mathcal S} s_j\;} \]
$\lambda$ existe et est unique dès que $\lvert\mathcal S\rvert \ge \lceil 1/c\rceil = 10$.

\etape{8}{L'exécution}
Achat au close du lendemain, $e = t+1$ ; \textbf{parts figées} entre deux coupes.
\[ q_i = V_e\,w_i / P_i(e), \qquad V_u = \textstyle\sum_i q_i P_i(u),\ \ u \in (e, e'] \]

\etape{9}{Le frottement}
\[ \mathrm{TO}_t = \tfrac12\textstyle\sum_i \bigl\lvert w^{\text{réal}}_i - w^{\text{dér}}_i \bigr\rvert,
   \qquad V_e \leftarrow V_e\,\bigl(1 - 2\,\mathrm{TO}_t\,\kappa\bigr) \]
$\kappa = 10$ points de base ; le facteur 2 parce qu'on paie à l'achat \textbf{et} à la vente.

\etape{10}{Les mesures}
\[ x_Y = \textstyle\prod_{t\in Y}(1+r_t) - \prod_{t\in Y}(1+r^{m}_t),
   \qquad t = \bar x \big/ (\sigma_x/\sqrt n) \]
\[ r - r_f = \alpha_1 + \beta\,(r^{m} - r_f) + \varepsilon \]
{\footnotesize\color{grisf} Facteurs quotidiens tirés de la
\href{https://mba.tuck.dartmouth.edu/pages/faculty/ken.french/data_library.html}{\color{bleu}Kenneth
R.\ French Data Library} --- séries {\urlstyle{tt}\url{F-F_Research_Data_Factors_daily}} et
{\urlstyle{tt}\url{F-F_Momentum_Factor_daily}}. Modèle de
\href{https://doi.org/10.1016/0304-405X(93)90023-5}{\color{bleu}Fama \& French (1993)}, augmenté du
momentum de \href{https://doi.org/10.1111/j.1540-6261.1997.tb03808.x}{\color{bleu}Carhart (1997)}.}\par\vspace{0.5mm}
\[ r - r_f = \alpha_4 + \beta_M (r^{m}-r_f) + \beta_S \mathrm{SMB} + \beta_H \mathrm{HML} + \beta_U \mathrm{MOM} + \varepsilon \]
\vspace{-1.2mm}
{\scriptsize
\begin{tabular}{@{}l@{~~}>{\raggedright\arraybackslash}p{1.65cm}>{\raggedright\arraybackslash}p{2.2cm}>{\raggedright\arraybackslash}p{2.2cm}@{}}
\toprule
 & le pari mesuré & $\beta > 0$ & $\beta < 0$ \\
\midrule
$\beta_S$ SMB & petites $-$ grosses & on ressemble aux petites & on ressemble aux grosses \\
$\beta_H$ HML & décotées $-$ chères & \emph{value} : banques, énergie & \emph{croissance} : la tech \\
$\beta_U$ MOM & gagnants $-$ perdants & on suit la tendance & on va à contre-courant \\
\bottomrule
\end{tabular}}
\par\vspace{0.5mm}
{\scriptsize\color{grisf} \textbf{S}mall \textbf{M}inus \textbf{B}ig : la capitalisation boursière
$\cdot$ \textbf{H}igh \textbf{M}inus \textbf{L}ow : le rapport valeur comptable $\div$ prix de marché
$\cdot$ \textbf{Mom}entum : les 12 derniers mois.}
\par\vspace{0.9mm}
$x_Y$ = excès d'une \textbf{année civile} ($n = 13$) ; les $\alpha$ se lisent sur les rendements journaliers
et sont annualisés ($\times 252$) ; intervalle de Student à 95~\%.


% ══════════════════════════════════════════════════════════════════════
\mouv{II --- Ce que ça donne}{2014--2026, brut de frais}

{\centering\scriptsize\setlength{\tabcolsep}{3.4pt}
\begin{tabular}{@{}l r c r r r@{}}
\toprule
 & excès/an \emph{(t)} & IC 95~\% & $\alpha_1$ \emph{(t)} & $\beta$ & NAV \\
\midrule
\textbf{M3 --- NANC} & $\mathbf{+3{,}40}$ \emph{(1,61)} & $[-1{,}2;+8{,}0]$
  & $+1{,}66$ \emph{(1,36)} & 1,080 & \textbf{662} \\
\textbf{M3 --- GOP} & $\mathbf{+1{,}24}$ \emph{(1,61)} & $[-0{,}4;+2{,}9]$
  & $+1{,}09$ \emph{(1,07)} & 1,004 & \textbf{574} \\
\midrule
\emph{repère : le SPY} & --- & --- & --- & \emph{0,980} & \emph{494} \\
\emph{repère : équipondéré} & \emph{$-1{,}38$ ($-1{,}45$)} & \emph{$[-3{,}5;+0{,}7]$}
  & \emph{$-0{,}95$ ($-0{,}72$)} & \emph{0,98} & \emph{430} \\
\emph{repère : $m_i$ seuls} & \emph{$-0{,}65$ ($-0{,}91$)} & \emph{$[-2{,}2;+0{,}9]$}
  & \emph{$-0{,}41$ ($-0{,}38$)} & \emph{0,99} & \emph{464} \\
\bottomrule
\end{tabular}\par}
\vspace{0.3mm}
{\scriptsize\color{grisf} rendement \emph{total} : NANC 16,6~\%/an $\cdot$ GOP 15,2 $\cdot$ SPY 13,9 ;
\textbf{9 années gagnantes sur 13} des deux côtés. Les repères tournent sur le \emph{même univers} --- seule
la pondération change. \textbf{Les deux intervalles contiennent zéro.}}\par

\fig{m3_nav.png}{à gauche la valeur du portefeuille, base 100 en février 2014, échelle log ; à droite
l'écart au SPY année par année (NANC $-11{,}6$ pt en 2022, $+19{,}7$ en 2023).}

\et{a}{D'où vient l'excès : la pondération, à univers identique}

{\centering\scriptsize
\begin{tabular}{@{}l r r r@{\hspace{2.5mm}} r r r@{}}
\toprule
 & \multicolumn{3}{c}{NANC} & \multicolumn{3}{c}{GOP} \\
\cmidrule(lr){2-4}\cmidrule(lr){5-7}
$w\propto$ & excès & $\alpha_1$ & NAV & excès & $\alpha_1$ & NAV \\
\midrule
$1$ \emph{(équipondéré)} & $-1{,}38$ & $-0{,}95$ & 430 & $-1{,}08$ & $-0{,}67$ & 448 \\
$m_i$ \emph{(élus seuls)} & $-0{,}65$ & $-0{,}41$ & 464 & $-0{,}78$ & $-0{,}17$ & 466 \\
$B_i$ \emph{(dollars purs)} & $+1{,}43$ & $+0{,}37$ & 554 & $+0{,}49$ & $+0{,}17$ & 523 \\
$\mathbf{B_i m_i}$ \emph{(M3)} & $\mathbf{+3{,}40}$ & $\mathbf{+1{,}66}$ & \textbf{662}
  & $\mathbf{+1{,}24}$ & $\mathbf{+1{,}09}$ & \textbf{574} \\
\bottomrule
\end{tabular}\par}
\vspace{0.3mm}
{\scriptsize\color{grisf} \textbf{le produit dépasse ses deux composantes} : ni les dollars seuls ni les élus
seuls ne rendent ce que rend leur produit. Le $\beta$ monte avec l'excès (0,98 $\to$ 1,08) --- c'est
l'$\alpha_1$, qui retire le marché, qui montre que l'essentiel ne vient pas de là.}\par

\et{b}{Le risque : est-ce un pari de style ?}

{\centering\scriptsize
\begin{tabular}{@{}l r r r r r r@{}}
\toprule
 & $\alpha_4$/an & $t$ & Mkt & SMB & HML & Mom \\
\midrule
M3 --- NANC & $\mathbf{+2{,}09}$~\% & \textbf{2,08} & 1,047 & $-0{,}131$ & $-0{,}146$ & $-0{,}076$ \\
M3 --- GOP & $+1{,}65$~\% & 1,79 & 0,979 & $-0{,}081$ & $+0{,}049$ & $-0{,}073$ \\
\emph{SPY (repère)} & \emph{$+0{,}04$} & \emph{0,12} & \emph{0,980} & \emph{$-0{,}123$} & \emph{$+0{,}012$} & \emph{$-0{,}009$} \\
\bottomrule
\end{tabular}\par}
\vspace{0.3mm}
\begin{pts}
\item $\alpha_4 > \alpha_1$ ($+1{,}66 \to +2{,}09$) : les facteurs de style n'expliquent pas l'excès, ils en
      \textbf{masquaient} une partie.
\end{pts}

\et{c}{Le résultat tient-il année après année ?}
Les chiffres ci-dessus écrasent douze ans en un nombre. On rejoue la régression sur une fenêtre
\textbf{glissante d'un an} (252 jours).

{\centering\scriptsize
\begin{tabular}{@{}l r c@{\hspace{3mm}} r c r@{}}
\toprule
 & \multicolumn{2}{c}{$\beta$ glissant} & \multicolumn{3}{c}{$\alpha$ glissant, \%/an} \\
\cmidrule(lr){2-3}\cmidrule(lr){4-6}
 & méd. & min--max & méd. & min--max & $>0$ \\
\midrule
NANC & 1,06 & 0,95--1,27 & $+2{,}01$ & $-11{,}0$ / $+13{,}2$ & \textbf{71~\%} \\
GOP & 1,00 & 0,90--1,13 & $+1{,}67$ & $-6{,}0$ / $+8{,}0$ & \textbf{71~\%} \\
\bottomrule
\end{tabular}\par}
\vspace{0.3mm}
{\scriptsize\color{grisf} \textbf{le même 71~\% des deux côtés} : l'$\alpha$ tient au-dessus de zéro 71~\% du
temps côté démocrate \emph{et} républicain, et sa médiane glissante ($+2{,}01$ / $+1{,}67$) est du même ordre
que l'$\alpha_1$ plein-échantillon ($+1{,}66$ / $+1{,}09$). \textbf{Réserve} : les fenêtres se chevauchent ---
ces chiffres \emph{décrivent}, ils ne testent pas.}\par

\et{d}{Le coût, et faut-il ralentir}

{\centering\scriptsize
\begin{tabular}{@{}l r r@{\hspace{4mm}} r r r r@{}}
\toprule
 & \multicolumn{2}{c}{rotation \%/an} & brut & net & coût & NAV nette \\
\cmidrule(lr){2-3}
 & médiane & \textbf{moy.} & & & & \\
\midrule
NANC & 65,6 & \textbf{71,8} & $+3{,}40$ & $+3{,}23$ & 0,17 & 649 \\
GOP & 69,2 & \textbf{76,9} & $+1{,}24$ & $+1{,}06$ & 0,18 & 562 \\
\bottomrule
\end{tabular}\par}
\vspace{0.3mm}
{\scriptsize\color{grisf} \textbf{c'est la moyenne qui fixe la facture}, pas la médiane. Le fonds réel publie
44, 62 puis 10~\% pour NANC.}\par
\vspace{0.4mm}
\begingroup\makeatletter\@beginparpenalty=10000\makeatother
{\color{bleu}\bfseries Le frein --- non retenu}\nopagebreak\par\nopagebreak\vspace{0.2mm}\nopagebreak
\begin{pts}
\item \textbf{ce qu'il fait} --- si rejoindre la nouvelle cible exige de tourner plus de $\tau_{\max}$, on
      ne parcourt qu'une fraction du chemin : $\lambda_R = \min\bigl(1,\ \tau_{\max}/\mathrm{TO}_t\bigr)$ ;
\item \textbf{à 25~\%/mois il est quasi inerte} --- M3 tourne 4,9~\% (NANC) et 5,3~\% (GOP) par mois en
      médiane, cinq fois sous le plafond : il ne mord que sur \textbf{1~\% des coupes} ;
\item \textbf{le balayage entier}, gain du parti le moins bien servi : $-0{,}04$ à 25~\%, $-0{,}05$ à
      10~\%, $\mathbf{+0{,}03}$ à 5~\%, $-0{,}48$ à 2~\%.
\end{pts}\endgroup

\et{e}{Le filtre anti-bruit $\theta$}

Un élu qui déclare 300 opérations le même jour ne prend pas 300 décisions : son courtier rééquilibre un
mandat. Le filtre est massif --- il retire \textbf{54~\%} du \$ démocrate et \textbf{58~\%} du républicain.

{\centering\scriptsize
\begin{tabular}{@{}l r r r r r@{}}
\toprule
 & $\theta=10$ & 20 & \textbf{50} \emph{(M4)} & 100 & $\infty$ \emph{(M3)} \\
\midrule
NANC --- excès & $+4{,}71$ & $+4{,}62$ & $\mathbf{+5{,}14}$ & $+4{,}75$ & $+3{,}40$ \\
\emph{\quad$\alpha_1$} & \emph{$+0{,}60$} & \emph{$+0{,}78$} & \emph{$+1{,}42$} & \emph{$+1{,}36$} & \emph{$\mathbf{+1{,}66}$} \\
GOP --- excès & $+3{,}02$ & $+2{,}14$ & $\mathbf{+1{,}44}$ & $+1{,}64$ & $+1{,}24$ \\
\emph{\quad$\alpha_1$} & \emph{$+2{,}77$} & \emph{$+1{,}88$} & \emph{$+1{,}54$} & \emph{$+1{,}60$} & \emph{$+1{,}09$} \\
\bottomrule
\end{tabular}\par}
\vspace{0.3mm}
{\scriptsize\color{grisf} balayage publié en entier, verdict lu à la valeur \emph{déclarée d'avance} ---
et aucune ne domine : GOP préfère 10, NANC préfère 50.}\par

\et{f}{Le netting des ventes --- testé, non retenu}
Le fonds réel nette : une vente réduit la position. On l'applique, \textbf{et rien d'autre} --- même fenêtre,
même score, même top-150, même plafond. De chaque achat on retire la part des ventes portant sur un lot de
moins de trois ans (FIFO $\gamma_3$) ; vendre un lot plus ancien ne change rien, comme au prospectus.

{\centering\scriptsize
\begin{tabular}{@{}l r r r r r@{}}
\toprule
 & M3 & \emph{avec netting} & écart & $t$ & $\alpha_1$ \\
\midrule
NANC & $+3{,}40$ & $+1{,}31$ & $\mathbf{-2{,}09}$ & 0,76 & $-0{,}30$ \\
GOP & $+1{,}24$ & $+1{,}06$ & $-0{,}18$ & 1,17 & $+0{,}90$ \\
\bottomrule
\end{tabular}\par}
\vspace{0.3mm}
{\scriptsize\color{grisf} \textbf{les deux partis se dégradent} : la règle qui a écarté le frein tranche donc
ici aussi. Mais l'ampleur n'a rien à voir --- $-2{,}09$ pt côté démocrate contre $-0{,}18$ côté républicain.
\textbf{Contrôle} : en neutralisant $\gamma_3$, la fonction redonne M3 à $10^{-12}$, ce qui garantit que
\emph{seul} le netting change.}\par

% ══════════════════════════════════════════════════════════════════════
\mouv{III --- Version 2}{un portefeuille qui n'achète que des ETF}

\textbf{M3 tourne en entier, inchangée} --- les dix étapes ci-dessus, à l'identique. \textbf{Trois étapes de
plus}, et elles viennent toutes \emph{après}.\par
\vspace{0.3mm}
{\footnotesize\color{grisf} Le livrable : \textbf{onze lignes} d'ETF, $+3{,}41$~\%/an net, écart au marché
$4{,}2$~\%.}\par

\etape{11}{L'instrument}
Chaque ligne devient l'ETF de son secteur ; les poids se somment, puis se renormalisent sur les seuls
secteurs cotés :
\[ w^{S} = \!\!\!\sum_{i\,:\,\mathrm{etf}(i)=S}\!\!\! w_i \ , \qquad
   w^{S} \leftarrow w^{S} \Big/ \textstyle\sum_{S' \text{ coté}} w^{S'} \]
\begin{pts}
\item le rattachement titre $\to$ GICS $\to$ SPDR est un \textbf{intrant} du dépôt, renseigné sur 100~\% des lignes ;
\item la clé est le ticker \emph{après} fusion --- celle de l'étape 5 ;
\item cinq tickers portent deux secteurs (BALL, CPAY, SONY, TFC, XYZ) : on retient le majoritaire ;
\item \textbf{la carte est datée.} XLRE naît en 10/2015 et XLC en 06/2018, mais leurs \textbf{indices} ne
      basculent qu'au 16/09/2016 et au 21/09/2018. Avant la bascule, chaque titre est rangé dans l'ETF de
      \textbf{son secteur d'alors} :
\end{pts}
\vspace{-0.6mm}
{\centering\scriptsize\setlength{\tabcolsep}{4pt}
\begin{tabular}{@{}l c l@{}}
\toprule
un titre aujourd'hui classé\dots & avant & à partir de \\
\midrule
immobilier & \textbf{XLF} & \textbf{XLRE}, le 16/09/2016 \\
communication, \emph{venue de la tech} & \textbf{XLK} & \textbf{XLC}, le 21/09/2018 \\
communication, \emph{venue de la conso} & \textbf{XLY} & \textbf{XLC}, le 21/09/2018 \\
\bottomrule
\end{tabular}\par}
\vspace{0.3mm}
{\scriptsize\color{grisf} \emph{venue de la tech} : Google, Meta, AT\&T, les jeux vidéo $\cdot$ \emph{de la
conso} : Netflix, Disney, les médias. Sans cette carte, un titre retenu par le score mais dont le secteur n'a
pas encore d'ETF serait \textbf{perdu} --- 11 titres en médiane, jusqu'à 20,8~\% du portefeuille.}\par

\etape{12}{Le signal : un seul portefeuille}
Le score ne compte plus les élus d'un camp, mais ceux \textbf{des deux} :
\[ m_i = \#\bigl\{\, m(k) \ :\ k \in \mathcal K^{\text{D}}(t) \cup \mathcal K^{\text{R}}(t) \,\bigr\} \]
Trois démocrates \emph{et} trois républicains donnent $m_i = 6$, pas 3 --- et le score compte au carré
(étape 5), donc \textbf{le consensus bipartisan est récompensé}.
\begin{pts}
\item retenir un camp serait un choix fait \emph{après} avoir vu les résultats ; celui-ci n'en demande aucun.
      Dernière coupe : 39 acheteurs démocrates, 60 républicains, \textbf{100 fusionnés}.
\end{pts}

{\centering\scriptsize\setlength{\tabcolsep}{2.5pt}
\begin{tabular}{@{}l r r@{\hspace{2.5mm}} r r@{\hspace{2.5mm}} r r@{}}
\toprule
 & \multicolumn{2}{c}{\textbf{2014--2026}} & \multicolumn{2}{c}{\textbf{2019--2026}} & & \\
\cmidrule(lr){2-3}\cmidrule(lr){4-5}
 & excès/an \emph{(t)} & $\alpha_1$ & excès/an & $\alpha_1$ & $\beta$ & NAV \\
\midrule
\multicolumn{7}{@{}l}{\emph{sur les 150 titres}}\\
démocrates & $+3{,}40$ \emph{(1,61)} & $+1{,}66$ & $+4{,}16$ & $+1{,}62$ & 1,08 & 662 \\
républicains & $+1{,}24$ \emph{(1,61)} & $+1{,}09$ & $+2{,}28$ & $+2{,}21$ & 1,00 & 574 \\
\textbf{unique} & $\mathbf{+2{,}51}$ \emph{\textbf{(1,93)}} & $\mathbf{+1{,}48}$ & $+3{,}29$ & $\mathbf{+1{,}87}$ & 1,05 & \textbf{629} \\
\midrule
\multicolumn{7}{@{}l}{\emph{sur les 11 ETF sectoriels}}\\
démocrates & $+1{,}78$ \emph{(1,66)} & $+0{,}59$ & $+1{,}92$ & $+0{,}20$ & 1,05 & 586 \\
républicains & $+0{,}78$ \emph{(1,02)} & $+0{,}41$ & $+1{,}12$ & $+0{,}65$ & 1,02 & 552 \\
\textbf{unique} & $\mathbf{+1{,}45}$ \emph{\textbf{(2,05)}} & $\mathbf{+0{,}46}$ & $+1{,}70$ & $\mathbf{+0{,}32}$ & 1,04 & \textbf{577} \\
\bottomrule
\end{tabular}\par}
\vspace{0.2mm}
{\scriptsize\color{grisf} lignes ETF mesurées avec la \textbf{carte datée} · le portefeuille unique porte les
deux plus hauts $t$ du document (1,93 et \textbf{2,05}), \textbf{aucun n'atteint le seuil} de 2,18 qu'exigent
treize années --- annexe D.}\par

\etape{13}{La forme livrable : le marché, plus un tilt}
\[ w \;=\; \underbrace{w^{\text{marché}}}_{\text{gratuit}} \;+\; \lambda\,\underbrace{(w^{\text{signal}} - w^{\text{marché}})}_{\text{le pari}} \]
$\lambda = 0$ donne le marché, $\lambda = 1$ la version sectorielle simple. Les deux grandeurs qui décident
de $\lambda$ --- l'\textbf{écart au marché} et ce qu'une unité d'écart rapporte :
\[ \mathrm{TE} = \sigma\bigl(r_p - r^{m}\bigr)\sqrt{252} \ , \qquad
   \text{rdt/écart} = \bar x \big/ \mathrm{TE} \]
\textbf{Les deux vecteurs.} $w^{\text{signal}}$ = M3 traduite en secteurs, c'est l'étape 11.
$w^{\text{marché}}$ se reconstruit depuis les prix, les onze SPDR \textbf{partitionnant} le S\&P 500 :
\[ r^{\text{SPY}} = \beta_1 r^{\text{XLK}} + \beta_2 r^{\text{XLF}} + \dots + \beta_{11} r^{\text{XLU}}
   + \varepsilon \]
Sous $\sum \beta = 1$ et $\beta \ge 0$, les $\beta$ tiennent lieu de poids sectoriels ($R^2$ de 0,984 à
0,999), estimés sur \textbf{l'année précédente} --- aucun regard vers l'avenir. Après le tilt, tout poids
négatif est ramené à zéro (long seul).

{\centering\scriptsize\setlength{\tabcolsep}{2.5pt}
\begin{tabular}{@{}l r r r r r@{}}
\toprule
$\lambda$ & excès net/an \emph{(t)} & $\alpha_1$ \emph{(t)} & $\beta$ & \textbf{écart marché} & rdt/écart \\
\midrule
0 \emph{(le marché)} & $-0{,}33$ \emph{(--0,3)} & $-0{,}65$ \emph{(--0,8)} & 1,00 & 2,1~\% & --- \\
1 & $+1{,}73$ \emph{(1,36)} & $+0{,}42$ \emph{(0,45)} & 1,04 & 2,6~\% & 0,68 \\
2 & $\mathbf{+3{,}41}$ \emph{\textbf{(1,82)}} & $\mathbf{+1{,}29}$ \emph{(0,87)} & 1,09 & \textbf{4,2~\%} & \textbf{0,83} \\
3 & $+4{,}31$ \emph{(1,84)} & $+1{,}71$ \emph{(0,93)} & 1,10 & 5,2~\% & 0,85 \\
\bottomrule
\end{tabular}\par}
\vspace{0.2mm}
{\scriptsize\color{grisf} l'\textbf{écart au marché} = écart-type annualisé de la différence avec l'indice : à
4,2~\%, deux années sur trois l'écart au SPY tient entre $-4{,}2$ et $+4{,}2$~\%. \textbf{Fenêtre 2020--2026,
net de 10~pb --- ne se compare pas au tableau ci-dessus}, qui couvre 148 coupes depuis 2014, brut de frais.}\par
\vspace{0.5mm}

\fig{etf_livrable.png}{à gauche le portefeuille depuis 2020 (base 100, log) ; à droite ce que chaque niveau
d'écart au marché rapporte --- la courbe est concave.}

\et{a}{Ce qui contraindra avant l'écart au marché}
\begin{pts}
\item \textbf{aucun plafond sectoriel} : celui de 10~\% porte sur les \emph{titres} et \textbf{ne survit pas}
      à l'agrégation --- cinq titres tech à 6--10~\% font un secteur à 46,6~\% (dém.) et 42,2~\% (rép.) ;
\item \textbf{et le tilt concentre} : le plus gros secteur pèse 38,7~\% à $\lambda = 0$ \emph{(la
      concentration du marché lui-même, en 2025)}, 47,8~\% à $\lambda = 1$, \textbf{61,2~\% à $\lambda = 2$}
      --- $N_{\text{eff}}$ de 6,5 à 4,6.
\end{pts}

\et{b}{Deux pistes écartées}
\begin{pts}
\item \textbf{32 instruments plus fins} (semi-conducteurs, biotechnologie, défense, banques régionales) :
      \textbf{dégrade} --- $+0{,}31$ et $-0{,}60$~\%/an contre $+1{,}50$ et $+0{,}45$. C'est l'instrument, pas
      le signal : \textbf{les ETF fins ne battent le secteur qu'ils découpent que 6 fois sur 21}, $-1{,}6$~pt/an ;
\item \textbf{recomposer plus souvent} : sans objet --- le signal est une \textbf{inclinaison}, pas une
      rotation : 85~\% des écarts de poids sont permanents côté démocrate, 76~\% côté républicain, et 2,2~\%
      du portefeuille seulement change de main d'une coupe à l'autre.
\end{pts}

\end{multicols}

\clearpage

% ══════════════════════════════════════════════════════════════════════
\vspace{1.5mm}{\color{bleu}\hrule height 1pt}\vspace{0.8mm}
{\color{bleu}\bfseries Annexe A --- comment on l'a vérifié}\par
\vspace{0.6mm}
{\footnotesize
Chaque ligne est un contrôle exécuté par le notebook, avec le chiffre qu'il imprime.\par}
\vspace{0.4mm}
{\centering\scriptsize
\begin{tabular}{@{}l l l@{}}
\toprule
\textbf{ce qui est vérifié} & \textbf{comment} & \textbf{résultat} \\
\midrule
la chaîne de mesure ne fabrique rien & le SPY passé au même moulin & $\beta_{\text{MKT}}$ 0,980 $\cdot$ $\alpha_4$ $+0{,}04$~\% ($t$ 0,12) \\
le moteur fait ce qui est écrit & un mois recalculé à la main, hors moteur & NAV 106,442585 des deux côtés \\
le plafond est bien $\min(c,\lambda v)$ & $\lambda$ retrouvé par bissection indépendante & écart max $9{,}7\cdot10^{-17}$ \\
le coût facturé est celui de la formule & $\mathrm{NAV}_{\text{nette}}/\mathrm{NAV}_{\text{brute}}$ contre $\prod_t(1-2\,\mathrm{TO}_t\kappa)$ & écart $4{,}8\cdot10^{-15}$ (661,57 $\to$ 648,69) \\
l'extraction est fidèle au papier & rejouée sur 150 coupes, au plafond de 8,5~\% & \textbf{18 chiffres reproduits} (NAV 672,35 / 612,25) \\
la sélection mord & cardinal de $\mathcal S$ sur les 148 coupes & min 99 (NANC) et 112 (GOP), \textbf{médiane 150} \\
aucune coupe ne dégénère & $\lvert\mathcal S\rvert \ge \lceil 1/c\rceil$ partout & 0 coupe sous 10 titres \\
les ETF n'entrent pas dans le flux & flux recompté après chargement des 11 SPDR & 113\,369 opérations, inchangé \\
le passage aux ETF ne change que l'agrégation & carte identité (chaque titre $\to$ lui-même) & redonne M3 à $2{,}8\cdot10^{-17}$ \\
le poids d'un secteur est la somme de ses titres & conservation, avec la vraie carte & vérifié sur 296 couples \\
le repli sur secteur non coté est au prorata & proportions relatives des survivants & vérifié sur 104 couples \\
M3 sur les titres n'a pas bougé & non-régression après le bloc III & NAV 661,57 et 573,57 \\
\bottomrule
\end{tabular}\par}
\vspace{0.5mm}
{\footnotesize\color{grisf} \textbf{Ce que ces contrôles ne disent pas} : ils garantissent que le calcul fait
ce que le protocole décrit --- ils ne disent rien de la \emph{portée} des résultats, qui est en annexe D.}\par

% ══════════════════════════════════════════════════════════════════════
\vspace{1.5mm}{\color{bleu}\hrule height 1pt}\vspace{0.8mm}
{\color{bleu}\bfseries Annexe B --- ce qui fixe chaque constante, hors de toute performance}\par
\vspace{0.6mm}
{\centering\scriptsize
\begin{tabular}{@{}l l l@{}}
\toprule
\textbf{étape} & \textbf{valeur} & \textbf{ce qui la fixe --- jamais notre rendement} \\
\midrule
1 $\cdot$ montant retenu & milieu de la fourchette déclarée & la règle publiée du fonds \\
3 $\cdot$ recomposition & mensuelle & un signal à 3 ans ne se renouvelle qu'un trente-sixième par mois \\
3 $\cdot$ départ du backtest & 28/02/2014 & la 1\textsuperscript{re} coupe où le plafond a une solution ($\ge10$ titres) \\
4 $\cdot$ fenêtre du signal & 756 j ($\simeq$ 3 ans) & la règle publiée du fonds \\
4 $\cdot$ date prise en compte & la \textbf{divulgation} $\delta$ & la date de transaction n'est pas investissable \\
6 $\cdot$ univers & les $N = 150$ plus gros scores & le prospectus annonce 100 à 200 lignes \\
7 $\cdot$ plafond de position & $c = 10$~\% & la médiane du fonds réel sur 13 dates, \emph{et} la limite UCITS \\
9 $\cdot$ coût facturé & 10 points de base & l'ordre de grandeur d'un aller-retour sur méga-capitalisation \\
\emph{hors M3} $\cdot$ filtre anti-bruit & $\theta = 50$ opérations & la description d'un gérant, pas un réglage \\
\bottomrule
\end{tabular}\par}
\vspace{0.4mm}
{\footnotesize\color{grisf} Le plafond mérite un mot : c'est le seul paramètre absent du prospectus. Deux
sources indépendantes, dont aucune n'est notre rendement, donnent la même valeur --- la plus grosse ligne du
fonds réel vaut \textbf{10,3~\% en médiane} sur ses treize états déposés à la SEC (entre 8,0 et 12,7~\%), et
\textbf{10~\% est la limite UCITS} par émetteur. Il en faut un : \textbf{sans plafond, la plus grosse ligne
monte à 36~\% en médiane et jusqu'à 87~\%}.}\par

% ══════════════════════════════════════════════════════════════════════
\vspace{1.5mm}{\color{bleu}\hrule height 1pt}\vspace{0.8mm}
{\color{bleu}\bfseries Annexe C --- le score sectoriel n'est-il qu'un reflet de la taille des secteurs ?}\par
\vspace{0.6mm}
{\footnotesize
La technologie pèse un tiers du portefeuille sectoriel côté démocrate --- mais un tiers du S\&P 500 aussi.
L'exposition sectorielle du marché se reconstruit depuis les prix (les onze SPDR couvrent l'indice, $\beta \ge 0$
et $\sum\beta = 1$) : \emph{un proxy des poids, pas les poids eux-mêmes}.\par}
\vspace{0.4mm}
{\centering\scriptsize\setlength{\tabcolsep}{5pt}
\begin{tabular}{@{}l r r r r@{}}
\toprule
 & $\rho$ de rang, portefeuille contre marché & écart absolu moyen & \multicolumn{2}{c}{la technologie (XLK)} \\
\cmidrule(lr){4-5}
médiane sur 2019--2025 & & & marché & portefeuille \\
\midrule
NANC & 0,85 & 3,00 pt & 29,36~\% & 31,71~\% \\
GOP & 0,88 & 2,38 pt & 29,36~\% & 19,75~\% \\
\bottomrule
\end{tabular}\par}
\vspace{0.4mm}
{\footnotesize
\textbf{Le classement, oui ; les poids, non.} Le portefeuille range les onze secteurs à peu près comme le
marché les pèse, mais s'en écarte de deux à trois points par secteur --- et en sens \emph{opposés} sur le
secteur qui décide de tout : les démocrates surpondèrent la technologie (31,7 contre 29,4~\%), les
républicains la sous-pondèrent de dix points (19,8~\%). C'est précisément cet écart que la Version 2 achète,
en laissant le reste à un indice.\par
\vspace{0.4mm}
\textbf{Et ce n'est pas un pari sur un seul secteur} : le plus gros contributeur ne pèse que \textbf{36~\%} de
l'excès côté démocrate (la technologie) et \textbf{34~\%} côté républicain (l'énergie) --- deux moteurs
\emph{opposés}.\par}

% ══════════════════════════════════════════════════════════════════════
\vspace{1.5mm}{\color{bleu}\hrule height 1pt}\vspace{0.8mm}
{\color{bleu}\bfseries Annexe D --- ce que cette fiche n'établit pas}\par
\vspace{0.6mm}
{\footnotesize
\textbf{La significativité.} Aucun excès mesuré ici n'est significatif : treize années (sept pour la
Version 2) ne suffisent pas à trancher, et les intervalles de confiance contiennent zéro.
\textbf{Les frais.} Seuls 10 points de base par aller-retour, à l'étape 9 --- aucun impact de marché, aucune
contrainte de liquidité, aucune capacité chiffrée. \textbf{Le portefeuille.} Long seul, sans levier, sans
vente à découvert. \textbf{Les facteurs.} Ceux de Kenneth French, repris tels quels ; ils s'arrêtent au
30/04/2026, si bien que les régressions portent sur 3\,059 jours au lieu de l'échantillon complet.
\textbf{La fidélité.} Reproduire le fonds réel est une \emph{autre} question --- cette fiche mesure une
stratégie.\par
\vspace{0.5mm}
\textbf{Le calendrier des ETF sectoriels.} La carte est datée (étape 11), sur des dates de bascule d'indice
établies par sources primaires SEC : l'immobilier quitte \textbf{XLF} le 16/09/2016 (formulaire 497 du
02/09/2016), la communication quitte \textbf{XLK} ou \textbf{XLY} le 21/09/2018 (SAI du 485BPOS du 15/06/2018)
--- soit onze et trois mois \emph{après} la première cotation de XLRE et de XLC. Le discriminant entre XLK et
XLY est l'\textbf{industrie} du titre ; les 46 sociétés disparues, pour lesquelles aucune source de données ne
répond plus, reçoivent une attribution \textbf{nominative} déclarée dans le notebook avant tout backtest.
\textbf{Ce qui reste non attribué --- huit tickers, tous des fonds.} \textbf{EDR} d'abord : sa description dit
\og{}\emph{Education Realty Trust}\fg{}, un REIT, et son secteur Communication Services est simplement
\emph{faux}.
Puis sept ETF --- \textbf{FDN}, \textbf{IXP}, \textbf{PNQI}, \textbf{VOX} (classés communication) et
\textbf{VNQ}, \textbf{IYR}, \textbf{SCHH} (classés immobilier). Ils passent tous le filtre
\texttt{asset\_class == "stock"}, et \emph{cinq d'entre eux portent} \texttt{asset\_type == "Mutual Fund"}
\emph{sur la même ligne} : la table se contredit. Deux entrent réellement dans le top-150 --- PNQI sur
10 coupes, SCHH sur 4. Les réattribuer aurait fait acheter un secteur au prétexte d'un fonds mal étiqueté ;
ils gardent donc l'ancien comportement. \textbf{La vraie correction serait de les retirer du flux}, mais elle
déplacerait toutes les ancres du projet --- chantier du pipeline, consigné et non traité ici.\par}

\end{document}
